\head{第八版新核查：第0章19.7--19.10}
\source{EGA IV第一部分，原分册印刷第104--116页；上传PDF第103--115页。作者勘误据IV第四部分，印刷第345、347页；上传PDF第344、346页}
本轮接续第七版的（19.6.7），连续处理19.7的章首至（19.10.4）末尾。第0章§19由此全部收入工作译稿。本轮新增32个原著编号条目、原编号公式（19.9.8.1）及两幅交换图。原文明确省略的证明仍标明省略，未以译者补写冒称原文。

本版正文采用已经核实的修订陈述；原文错误及理由写入脚注，较长反例与补证保存在本报告。必须特别说明：（19.7.1.5）的较弱拓扑版本本轮\textbf{没有取得完整验证，也没有取得反例}。正文明确给出足以覆盖后文实际应用的完备局部充分版本，底本较弱版本的完整译文保存在下文H8-P01。该处理是范围限定，不是把一个尚未否定的原命题宣布错误。

\section*{本轮结论概览}
\addcontentsline{toc}{section}{本轮结论概览}
\begin{longtable}{@{}p{.17\linewidth}p{.19\linewidth}p{.55\linewidth}@{}}
\toprule 编号 & 原条目 & 判定与正文处理\\\midrule\endhead
H8-E01 & 19.10.3(i) & 一般离散环范围内的原断言有反例；改为核理想 $I=I^2$ 的准确判据，并给双向证明。\\
H8-C01--C03 & 19.8.3--19.8.8 & 作者已有正式勘误：局部同态、引文号、$B$-模。已写入正文。\\
H8-W01 & 19.8.7(iii) & 删除对已撤销19.3.12的有效性依赖；仅保留19.3.10支持的范围。\\
H8-T01--T03 & 19.7.1.6；19.8.8；19.10.3(iv) & 分别处理引用错号、非唯一提升误称唯一、零基环边界。\\
H8-P01--P02 & 19.7.1.5；19.7.3(i) & 较弱拓扑版本待核；正文采用已证明的充分版本，原文完整存档。\\
H8-N01--N04 & 各相关条目 & 参数记号、$W(k)$含义、历史问题状态及相对形式光滑的逻辑前提。不是四个错误定理。\\
\bottomrule
\end{longtable}

\audit{H8-E01}{满同态形式光滑的准确判据}{原断言在其所列的一般范围内不成立；有反例及完整修订证明}
\oldauditsrc{（$0_{\mathrm{IV}}$，19.10.3，i），原分册第116页／PDF第115页}
原文在证明满同态 $\rho:A\to B$ 形式非分歧后，又说：“反之，$B$ 只有在 $\rho$ 为双射时才形式光滑。”该段仅假设环离散，没有局部性、诺特性或核有限生成的限制，因此不能把局部诺特环中可能成立的结论移用到这里。

\lead{最小型反例。}
取非零域 $k$，$A=k\times k$，$B=k$，并令 $\rho(a,b)=a$。核为 $I=0\times k$，满足
\[
 I\ne0,\qquad I^2=I.
\]
因此 $\rho$ 满而不单。令 $e=(0,1)$；它既属于 $I$ 又满足 $e^2=e$。给定离散 $A$-代数 $E$、幂零理想 $J\subset E$ 和 $A$-同态 $B\to E/J$，结构同态 $\alpha:A\to E$ 将 $e$ 映到 $J$。$\alpha(e)$ 一方面幂零，另一方面幂等，故为零：若 $J^N=0$，则 $\alpha(e)=\alpha(e)^N=0$。所以 $\alpha$ 通过第一投影唯一分解，给出 $B\to E$。这不仅是形式光滑，而且是形式平展的反例。

\lead{准确版本。}
对任意满环同态 $\rho:A\to B$，设 $I=\Ker\rho$。以下条件等价：
\[
\begin{aligned}
 &B\text{ 是形式光滑 }A\text{-代数};\\
 &I=I^2;\\
 &B\text{ 是形式平展 }A\text{-代数}.
\end{aligned}
\]
\begin{proof}
满同态本来就形式非分歧，所以只须证明前两个条件等价。

假设形式光滑。以 $A/I^2$ 为测试 $A$-代数，取平方零理想 $I/I^2$。$B=A/I$ 的恒等同态提升为 $A$-同态 $s:B\to A/I^2$。对每个 $a\in I$，因 $s$ 为 $A$-同态，其在 $A/I^2$ 中的类等于 $s(0)=0$，故 $a\in I^2$。于是 $I\subseteq I^2$；反向包含恒成立，故 $I=I^2$。

反之，设 $I=I^2$。给定测试代数 $E$、幂零理想 $J$、商同态 $\pi:E\to E/J$ 及 $A$-同态 $f:B\to E/J$。令 $\alpha:A\to E$ 为结构同态。交换条件为 $\pi\alpha=f\rho$，所以 $\alpha(I)\subseteq J$。若 $J^N=0$，则由 $I=I^N$ 得 $\alpha(I)\subseteq J^N=0$。因此存在唯一环同态 $\widetilde f:A/I\to E$ 使 $\widetilde f\rho=\alpha$，并且 $\pi\widetilde f=f$。这正是所需唯一提升。
\end{proof}

\lead{与局部情形的界限。}
若再假设 $A$ 为局部环，$I$ 为有限生成真理想，则 $I\subseteq\mathfrak m_A$，$I=I^2$ 及Nakayama引理给出 $I=0$。所以在此附加范围内，原句可以成立。特别地，这不反驳19.7讨论的诺特局部环形式光滑性判据。一般判据本身只涉及提升性质，不能不加条件地据此宣称有限表示或平坦性。

\lead{外部交叉核验。}
Stacks Project的形式光滑定义及提升到平方核商的判据，与上述证明一致（Tag 00TH）；形式平展的定义为唯一提升（Tag 00UP）。反例也可看成幂等元局部化：取 $s=(1,0)$，$A_s\simeq k$，而局部化形式平展（Tag 00UP，Lemma 10.150.5）。这些仅用于交叉核验；正文中的判据及证明由本次核查明确给出，不冒称底本或作者勘误已有此修订。

\audit{H8-C01}{素局部环的唯一性必须连同局部同态理解}{作者正式勘误；不是本次新发现}
\oldauditsrc{19.8.3，第110页；19.8.4，第111页；勘误清单三Err$_{\mathrm{IV}}$，9，第347页}
正式勘误要求在19.8.3的两处及19.8.4的开头，将“同态”换为“局部同态”。对局部环 $(A,\mathfrak m)$，结构映射 $\mathbf Z\to A$ 的逆像 $p\mathbf Z$ 是素理想，唯一给出局部映射 $\mathbf Z_{p\mathbf Z}\to A$。若不要求局部性，例如 $A=\mathbf Q$，每个素数 $q$ 都给出保单位同态 $\mathbf Z_{q\mathbf Z}\to\mathbf Q$；不同素局部环不再被唯一选定。

正文在相应位置直接使用“局部同态”，脚注说明作者修订，书末保留正式勘误译文。此处不把一般环同态混同于局部环同态。

\audit{H8-C02}{非唯一性备注的引文号与平方零理想}{作者引文勘误；另排除识别误差}
\oldauditsrc{19.8.7，第111页；勘误清单三A类，第345页}
原文三次引用21.5.5，作者明确要求改为21.5.3。正文已作相应替换。原扫描页同时确认条件是 $\mathfrak j^2=0$；自动提取会漏掉上标 $2$，从而误成 $\mathfrak j=0$。后者当然不能作为提升不唯一的例子，因此本轮以页面图像为准，没有误报为作者逻辑错误。

作为独立检验，取 $K=\mathbf Q(t)$、$C=K[\varepsilon]/(\varepsilon^2)$、$\mathfrak j=(\varepsilon)$，令 $u:K\to C/\mathfrak j=K$ 为恒等。对任一 $a\in K$，
\[
 v_a(f)=f+\varepsilon a\frac{df}{dt}\qquad(f\in K)
\]
为提升 $u$ 的环同态。Leibniz法则及 $\varepsilon^2=0$ 保证乘法性；它保单位，且 $v_a(t)=t+a\varepsilon$。当 $a\ne b$ 时，两提升不同。所有非零 $f\in K$ 的像在 $C$ 中都是单位，所以 $v_a$ 是局部同态。此例正说明“同构类唯一”不等于“给定剩余域同态的提升唯一”。

\audit{H8-C03}{有限性证明中的基环}{作者正式排印勘误}
\oldauditsrc{19.8.8，第112页倒数第2行；勘误清单三A类，第345页}
原文在证明 $A$ 是 $B$ 的有限代数时，误写“所以 $A$ 是有限型 $A$-模”。正式勘误将其改为 $B$-模。正文采用正确表述。

推理链为：$B$ 完备、$\mathfrak nA$ 为 $A$ 的定义理想、剩余扩张有限，故 $A$ 是拟有限 $B$-模；再用（$0_{\mathrm I}$，7.4.1）得到 $A$ 对 $B$ 有限。其后才可结合 $\dim A=\dim B$ 及整性证明 $B\to A$ 单射。若只说 $A$ 作为自身上的模有限生成，不能支持这一步。

\audit{H8-W01}{不再用已撤销命题扩张Cohen提升范围}{失效引用的修订，不是否定Cohen定理}
\oldauditsrc{19.8.7(iii)，第111页；Err$_{\mathrm{IV}}$，8，第347页}
原文说，可借19.3.10和19.3.12减弱Cohen定理中对目标环 $C$ 的假设。然而19.3.12已由作者正式删除，理由是原证明使用的若干集合不是线性流形。第七版正文已经处理了该条自身的撤销状态；本轮进一步处理这里对它的依赖。

现在正文只保留19.3.10实际支持的结论：对于相容的拓扑基环结构，$C$ 可度量、分离完备，$\mathfrak j$ 闭，且 $\mathfrak j^n\to0$ 时，连续同态可以提升。并未因撤销19.3.12而声称所有其他推广均不成立；也没有把19.8.6的完备诺特局部环版本撤销。原引用差异均在脚注H8-W01中明示。

\audit{H8-T01}{商环完备性的引文}{引文定位修订}
\oldauditsrc{19.7.1.6，第108页}
原文用（$0_{\mathrm I}$，6.3.5）说明 $B/\mathfrak mB$ 完备。所需的是完备诺特局部环中理想闭、商环完备的性质，属于7.3的完备化讨论。正文改引7.3.5及商完备性质，并明确 $B$ 已在该步骤前被换成完备化。本项不改变定理的条件或结论。

\audit{H8-T02}{提升的选择不应误称唯一}{措辞修订}
\oldauditsrc{19.8.8证明，第112页}
证明开头由Cohen提升定理选择了 $u:W\to A$。后文却称它为“唯一同态”，与19.8.7明确指出的非唯一性不符。正文改为“所选同态”。后面的构造只使用该映射确实存在、局部、诱导指定的剩余域同构，并将整数 $p$ 映到 $p\cdot1$，不使用唯一性，所以结论不受影响。

\audit{H8-T03}{非空多项式代数的零基环边界}{边界条件补足}
\oldauditsrc{19.10.3(iv)，第116页}
若基环 $A=0$，则对任意变量集 $I$，$A[X_\alpha]_{\alpha\in I}=0$，结构同态为恒等，因而形式平展。故原文“$I\ne\varnothing$ 时不形式平展”须补 $A\ne0$。

在 $A\ne0$、$I\ne\varnothing$ 时，选 $\alpha_0\in I$，取测试代数 $E=A[\varepsilon]/(\varepsilon^2)$、理想 $J=(\varepsilon)$。两个 $A$-同态将所有其他变量都映为零，而分别将 $X_{\alpha_0}$ 映为 $0$ 和 $\varepsilon$。它们模 $J$ 相同，但本身不同。因此形式非分歧失败。正文直接列出非零基环条件，脚注给出例外和检验。

\audit{H8-P01}{同构提升引理的拓扑范围}{原较弱版本待核；给出足够用于后文的已验证版本}
\oldauditsrc{19.7.1.5，第108页}
\lead{底本原引理的完整译文。}
设 $A$ 为诺特环，$\mathfrak j$ 为其理想，$B,B'$ 为诺特局部 $A$-代数，结构同态 $A\to B$、$A\to B'$ 对 $A$ 的 $\mathfrak j$-预进制拓扑连续。假设：$1^\circ$ $B,B'$ 对 $\mathfrak j$-预进制拓扑完备；$2^\circ$ $B$ 为形式光滑 $A$-代数；$3^\circ$ $B'$ 为平坦 $A$-模。置 $A_0=A/\mathfrak j$。若给定 $A_0$-同构 $u_0:B\otimes_AA_0\to B'\otimes_AA_0$，则存在 $A$-同构 $u:B\to B'$，满足 $u\otimes1=u_0$；因而 $B'$ 形式光滑且 $B$ 对 $A$ 平坦。

\lead{底本原证明的完整译文。}
置 $B_0=B\otimes_AA_0$、$B'_0=B'\otimes_AA_0$。若 $\mathfrak n,\mathfrak n'$ 为 $B,B'$ 的极大理想，则因 $\mathfrak jB\subseteq\mathfrak n$、$\mathfrak jB'\subseteq\mathfrak n'$，两环的 $\mathfrak j$-预进制拓扑分离。又因 $\mathfrak jB'$ 对 $\mathfrak j$-预进制拓扑闭，连续复合 $B\to B_0\xrightarrow{u_0}B'_0$ 根据（19.3.10）分解为 $B\xrightarrow{u}B'\to B'_0$，其中 $u$ 为连续 $A$-同态。显然 $u\otimes1=u_0$，由（19.7.1.4），$u$ 双射。

\lead{本轮为何不将较弱版当作已经核定。}
本小节章首约定：局部环默认取极大理想拓扑。原引理只明确 $\mathfrak j$-进制完备，却直接援引要求目标环在所用拓扑下完备的19.3.10。若将 $B'$ 改用 $\mathfrak jB'$-拓扑，还必须说明源 $B$ 的形式光滑性与连续性是相对于哪一种拓扑；不能把两种拓扑无条件互换。

这不是一个只有名词上的区别。取 $R=\mathbf Q[x]_{(x)}$，$B=R[[t]]$，$J=(t)$。$B$ 为诺特局部环，极大理想为 $(x,t)$，并且按定义对 $J$-进制拓扑完备。但 $B$ 不对 $(x,t)$-进制拓扑完备：否则商环 $B/(t)=R$ 也应对 $(x)$-进制拓扑完备，而 $R$ 的完备化为 $\mathbf Q[[x]]$，严格大于 $R$。例如 $\sqrt{1+x}$ 的二项式部分和在 $(x)$-进制拓扑中是柯西列；它在 $R$ 中没有极限，因为若 $r(x)\in\mathbf Q(x)$ 满足 $r(x)^2=1+x$，则在不可约多项式 $1+x$ 处，左边的赋值为偶数，右边却为1。

此例\textbf{仅说明两种完备性不同，不是原引理的反例}。原文可能隐含另一拓扑约定，或存在补充提升论证；本轮尚未核定。因此不将其列入已反驳断言，而将正文限制为可直接验证且足以支持后文的版本。

\lead{已核实的充分版本及证明。}
在原引理中加强为：$B,B'$ 均对各自极大理想拓扑完备，形式光滑性也按章首的局部环拓扑理解。结构同态的连续性仍要求 $\mathfrak jB\subseteq\mathfrak n$、$\mathfrak jB'\subseteq\mathfrak n'$。这样两环也 $\mathfrak j$-进制分离完备。证明如下：一个 $\mathfrak j$-进制柯西列自动是极大理想进制柯西列；其极限存在，而且对每个 $r$，理想 $\mathfrak j^rB$ 或 $\mathfrak j^rB'$ 都在极大理想拓扑中闭，故极限满足全部模 $\mathfrak j^r$ 条件。分离性则由 $\bigcap_r\mathfrak n^r=0$ 得到。

局部环同构 $u_0$ 对极大理想拓扑连续；目标 $B'$ 对其极大理想拓扑完备，$\mathfrak jB'$ 闭且其幂趋于零。由19.3.10可提升得到 $u:B\to B'$。再以 $N=B$、$M=B'$ 应用19.7.1.4，利用 $B'$ 对 $A$ 平坦及模 $\mathfrak j$ 的同构性，得 $u$ 双射。

\lead{依赖审计。}
19.7.1.6先把 $A,B$ 换成各自的完备化，随后构造的 $B'$ 也是完备诺特局部环，因此实际应用落在上述充分版本内。19.7.2中 $B_0$、$B$、$B'$ 及它们的相关商同样是完备诺特局部环。因此本版并未借待核的较弱拓扑范围来完成这些主定理的证明。此点也与Stacks Project Tag 07NQ的诺特局部环形式光滑性与平坦性判据相符；该交叉核验不用于宣布较弱版本错误。

\audit{H8-P02}{19.7.3的弱完备性备注与历史问题}{继承H8-P01的范围限定；不作现代研究状态断言}
\oldauditsrc{19.7.3(i)，第109页}
原文首句为：“应注意，在19.7.2中，只假设 $B,B'$ 对 $\mathfrak j$-预进制拓扑完备时，唯一性断言仍成立。”这依赖19.7.1.5的较弱版本，故在正文脚注中完整保留，但未作为已核定推广用于后续证明。

后面的存在性问题和在特定情况下对模 $p^n$ 平坦提升的改写，均按原文译出。“我们不知道”改用明确历史指向的“原著接着说，尚不知道”，不冒充2026年的文献调查结论。本轮没有研究或解决该历史问题，也没有以其曾经未解决为理由停止当前翻译。

\audit{H8-N04}{相对形式光滑的逻辑与记号}{防止翻译改变原概念，不是原命题错误}
\oldauditsrc{19.9.1--19.9.8，第114--115页}
原扫描中的 $\Lambda$ 与 $A$ 非常接近，文本提取多处将二者都识别成 $A$。本轮以原页逐项核对，保留三个层级
\[
 \Lambda\longrightarrow A\longrightarrow B.
\]
相对形式光滑要求的是：给定连续 $A$-同态 $B\to C/J$，\emph{若已有连续 $\Lambda$-同态提升到 $C$}，就存在连续 $A$-同态提升。不能删除斜体前提，否则会把“相对形式光滑”误译为通常形式光滑。

特别地，19.9.2(i)确为“$A$ 是相对于 $\Lambda$ 形式光滑的 $\Lambda$-代数”，而非声称任意 $\Lambda$-代数通常形式光滑。19.9.3的结论也是 $B\otimes_\Lambda A'$ 作为\emph{$B$-代数}相对于 $\Lambda$ 形式光滑。19.9.4的四个版本分别改变 $B$、$A$、$\Lambda$ 的完备化，所有帽号均按图像核对。19.9.8的扩张群为
\[
 \Exalcotop_{A/\Lambda}(B,L)
 =\Ker\bigl(\Exalcotop_A(B,L)\to\Exalcotop_\Lambda(B,L)\bigr),
\]
离散版本则为 $\Exalcom_{A/\Lambda}(B,L)$，不能把相对下标漏成 $A$。

\section*{新增条目覆盖与未译范围}
\addcontentsline{toc}{section}{新增条目覆盖与未译范围}
19.7收录3个主要条目及19.7.1.1--19.7.1.6六个子编号，共9个；19.8收录19.8.1--19.8.11，共11个；19.9共8个；19.10共4个，合计32个。条目计数仅检查覆盖，不是正确性认证。

当前第0章连续译至19.10.4，即§19末尾。下一接续为§20“导子与微分”的章首及20.1.1；本轮没有把已看到但未译入的§20章首计入完成范围。第0章§§20--23、第IV章主体及第V卷预备笔记余文仍未译完；全部旧译文也没有在本轮重新逐句校勘。

\section*{外部交叉核验的来源界限}
\addcontentsline{toc}{section}{外部交叉核验的来源界限}
本轮翻译底本为用户上传的法文原版及原作者的勘误，不用网页提要替代原文。下列现代资料仅用于核查相应定义与已知主定理，未将其内容冒充原著：Stacks Project，Tag 00TH（形式光滑映射）；Tag 00UP（形式平展映射及局部化）；Tag 07EA（拓扑环的形式光滑性）；Tag 07NQ（诺特局部环的形式光滑性、平坦性与特殊纤维）。查询日期：2026年11月15日。标签是稳定检索标识，网页条目序号可能随更新变化。
